\section{Implemented Pipeline}

This section describes the concrete algorithmic design used by the current code. Unlike the theory section, which derives replay estimators and the local Fisher geometry, the present section focuses on operational choices: readiness checks, probe growth policy, production collection, numerical stabilization, and acceptance logic. At the software level, the current sampling and replay machinery is built on Molly, and the reverse-mode parameter derivatives are obtained with Enzyme \citep{greener2024molly,moses2020enzyme}.

The current code organizes the method as a macro-epoch loop. At the beginning of each macro epoch, the pipeline constructs one AWH simulation per thermodynamic leg with the current physical parameters. Restart coordinates and velocities may be reused from the previous macro epoch. The AWH bias is reused only when restart continuity is preserved; a cold restart from the input structure does not reuse a previous bias profile. For each leg, the macro epoch follows the control flow summarized below.

The optimizer never consumes data from an actively adapting AWH bias. The following summary is the only authoritative end-to-end control-flow description in this document; the later subsections expand individual stages rather than introducing a second competing summary.

\subsection{Authoritative Control-Flow Summary}
\label{sec:authoritative_control_flow}

\begin{enumerate}
    \item \textbf{Stage A.} Advance each adapting AWH leg in blocks until the cheap readiness checks pass for the required stable streak.
    \item \textbf{Stage B.} Clone the current leg into a frozen-bias probe, select probe frames, rebuild replay templates, and require split convergence together with support-aware replay/AWH parity.
    \item \textbf{Frozen production.} Once every leg passes Stage B, launch frozen-bias production segments and apply the production frame-selection rules before any estimator is assembled.
    \item \textbf{Offline replay.} Evaluate the selected frames under the current Hamiltonians to assemble endpoint free energies, state observables, latent-space gradients, candidate-ensemble weights, and ESS diagnostics.
    \item \textbf{Trust-region base step.} Build the empirical Fisher metric, apply any reference-penalty diagonal regularization, use Jacobi preconditioning and eigentruncation, and scale the base step to the KL target.
    \item \textbf{Line search and acceptance.} Backtrack on the proposed update until candidate-ensemble ESS, per-pool drift bounds, and the tolerance-normalized objective test all pass; otherwise keep the best accepted state and return to fresh resimulation.
\end{enumerate}

\subsection{Stage A: Cheap Readiness Screening}

Stage A is a low-cost check applied after each AWH block. It does not attempt to estimate free energies directly. Instead, it asks whether the adaptive AWH run is mixed and statistically mature enough to justify a frozen-bias probe. The current implementation monitors:

\begin{itemize}
    \item recent mean bias-change magnitude,
    \item lambda-history effective sample size,
    \item Molly's linear-stage effective sample count $N_{\mathrm{eff}}$,
    \item full round trips between the coupled and decoupled states,
    \item recent endpoint occupancy,
    \item solvent-tail occupancy floors for late Lennard-Jones states in the
    solvent leg.
\end{itemize}

Only after a leg passes these Stage A criteria for a required stable streak does the pipeline launch a Stage B probe. Repeated Stage B failures can increase the required stable streak or introduce cooldown periods before the next probe.

\subsection{Stage B: Frozen-Bias Probe and Parity Validation}

Stage B is the decisive readiness check. The current leg is cloned into a frozen reference simulation, the AWH bias is held fixed, bias-update accumulators are cleared, and a short probe trajectory is generated. The probe frames are then processed offline:

\begin{enumerate}
    \item discard an initial fraction of the probe,
    \item apply stride and minimum/maximum retained-frame rules,
    \item rebuild CPU replay templates for all lambda states,
    \item evaluate each retained frame under every state Hamiltonian.
\end{enumerate}

The resulting energy matrix is used to compute two independent Stage B diagnostics.

\paragraph{Split-half convergence.}

The retained frames are split into temporal halves. A coupled-minus-decoupled free energy is computed independently from each half, and the split gap is the absolute difference between those two estimates.

\paragraph{Replay/AWH parity.} 

The same retained frames are used to reconstruct a full lambda free-energy profile from the frozen reference mixture. That replayed profile is aligned to the coupled endpoint and compared against the frozen AWH profile that generated the frames. The implementation tracks:

\begin{itemize}
    \item the raw parity gap over all states,
    \item the supported parity gap over states whose reweighting ESS exceeds a     configured support threshold,
    \item a separate endpoint parity gap,
    \item minimum support coverage over the full lambda schedule.
\end{itemize}

The default gate is therefore \emph{support-aware} but not support-blind: a leg must satisfy split convergence, supported parity, endpoint parity, and support coverage simultaneously.

\paragraph{Split-only continuation.}

The current code distinguishes true split-only failures from genuine parity failures. If split convergence fails while supported parity, endpoint parity, and  support coverage all still pass, the pipeline does \emph{not} return immediately to adaptive AWH sampling. Instead, it continues the same frozen probe in place, appending one more segment under the unchanged frozen bias. The raw logger history is concatenated and the entire frame-selection procedure is rerun on the combined probe so that discard, stride, and frame-cap rules remain consistent. This continuation ends as soon as the failure mode is no longer split-only or the configured segment budget is exhausted.

\paragraph{Retry policy.}

When Stage B fails for lack of support or split convergence, the next probe is grown. When it fails on parity, the probe length is held fixed and the leg is  sent back to Stage A so the adaptive bias can continue to improve. Near-pass behavior is leg-specific: by default, the solvent leg keeps the same probe length on a near-pass rather than shrinking it, while the vacuum leg may still shrink.

\subsection{Frozen-Bias Production Artifacts}

Once every leg in the cycle has passed Stage B, the pipeline launches a second frozen-bias simulation for production collection. For each leg, the current code stores:

\begin{itemize}
    \item the logged coordinates, active lambda history, and, where relevant, volume history,
    \item the frozen bias metadata,
    \item precomputed replay neighbor data,
    \item an ensemble-evaluation cache for CPU replay,
    \item the reference energy matrix $u_{\mathrm{ref},\ell}$ over the selected production frames and all lambda states.
\end{itemize}

These production artifacts are the only data used by the optimizer. Before replay begins, the raw frozen-bias logger is passed through a production-frame selection rule that mirrors the Stage B probe treatment: stride thinning, minimum-frame backfilling, evenly spaced maximum-frame capping, and optional leading-frame discard are all applied under the unchanged frozen bias. Every replay quantity used by optimization---including $u_{\mathrm{ref},\ell}$, endpoint free energies, observable estimates, observable gradients, and ESS diagnostics---is built on this selected frame set rather than on the denser raw logger.

The production duration can also differ by leg type. In the bundled defaults, solvent-like legs use a longer frozen production segment than vacuum legs because low-frequency observables such as density are typically the noisiest optimization targets. The implementation can furthermore append multiple back-to-back production segments under the same frozen bias, which is valid because the reference mixture remains unchanged throughout the append.

\subsection{Configuring User Targets and Observables}

The pipeline is intended to be configured directly from user code rather than
only through the bundled hydration example. At the outermost level, a run is
defined by a thermodynamic cycle, a simulation configuration, an optimization
configuration, and an optional list of explicit training targets:

\begin{enumerate}
    \item define the legs in \texttt{ThermodynamicCycleConfig},
    \item choose the trainable parameter pools and numerical settings in \texttt{SimulationConfig} and \texttt{OptimizationConfig},
    \item populate \texttt{sim\_cfg.targets} with any mix of cycle free-energy targets and state-observable targets,
    \item call \texttt{run\_pipeline(; sim\_cfg=..., opt\_cfg=...)}.
\end{enumerate}

\paragraph{Default and explicit target lists.}

If \texttt{sim\_cfg.targets} is left as \texttt{nothing}, the code constructs a
single default \texttt{CycleFreeEnergyTarget} from
\texttt{cycle.target\_dG\_kcal\_mol}. Supplying an explicit target list replaces
that default and allows arbitrary mixtures of:

\begin{itemize}
    \item \texttt{CycleFreeEnergyTarget}, which matches the full
    coefficient-weighted thermodynamic-cycle free energy assembled from all
    configured legs,
    \item \texttt{StateObservableTarget}, which matches one replayed observable
    evaluated on one physical state of one named leg.
\end{itemize}

The following configuration pattern is the intended user-facing entry point:

\begin{verbatim}
cycle_cfg = AWHGrads.ThermodynamicCycleConfig(
    legs=[solvent_leg, vacuum_leg],
    include_standard_state_correction=true,
    target_dG_kcal_mol=-5.01,
)

targets = AWHGrads.AbstractTrainingTarget[
    AWHGrads.CycleFreeEnergyTarget(
        name=:hydration_free_energy,
        target_dG_kcal_mol=cycle_cfg.target_dG_kcal_mol,
        tolerance_kcal_mol=0.05,
        weight=1.0,
    ),
    AWHGrads.StateObservableTarget(
        name=:solvent_density,
        leg=:solvent,
        state=:coupled,
        observable=AWHGrads.MassDensityObservable(),
        target_value=0.99815,
        tolerance=0.02,
        weight=1.0,
        unit_label="g/mL",
    ),
]

sim_cfg = AWHGrads.simulation_config_with(
    AWHGrads.default_simulation_config();
    cycle=cycle_cfg,
    targets=targets,
    parameter_pools=parameter_pools,
)

opt_cfg = AWHGrads.default_optimization_config()
runtime = AWHGrads.run_pipeline(; sim_cfg=sim_cfg, opt_cfg=opt_cfg)
\end{verbatim}

\paragraph{What a state observable target means.}

A \texttt{StateObservableTarget} is always tied to:

\begin{itemize}
    \item one leg name, e.g. \texttt{:solvent},
    \item one physical state selector, either \texttt{:coupled},
    \texttt{:decoupled}, or an explicit state index,
    \item one observable object or callable,
    \item one scalar reference value, tolerance, weight, and display label.
\end{itemize}

Operationally, the pipeline replays each retained frozen-bias production frame
under the requested lambda state, rebuilds a CPU replay \texttt{Molly.System}
for that state, evaluates the observable on that replayed system, and then
reweights those per-frame values with the same frozen reference mixture used by
the free-energy estimators. This is why cycle free energies and arbitrary state
observables can appear in one joint objective: both are derived from the same
selected frozen-bias production frames, the same replay Hamiltonians, and the
same state-conditioned importance weights.

\paragraph{Simple custom observables.}

For the common case, a user does not need to modify the pipeline internals at
all. Any callable that accepts the replayed \texttt{Molly.System} and returns a
finite scalar \texttt{Real} can be used directly as the \texttt{observable}
field of a \texttt{StateObservableTarget}. A minimal example is:

\begin{verbatim}
struct SoluteChargeObservable
    atom_idxs::Vector{Int}
end

function (obs::SoluteChargeObservable)(sys)
    return sum(Molly.charges(sys)[i] for i in obs.atom_idxs)
end

targets = AWHGrads.AbstractTrainingTarget[
    AWHGrads.StateObservableTarget(
        name=:solute_charge_sum,
        leg=:solvent,
        state=:coupled,
        observable=SoluteChargeObservable(collect(1:9)),
        target_value=0.0,
        tolerance=0.01,
        unit_label="e",
    ),
]
\end{verbatim}

The implementation evaluates such observables frame-by-frame on the replayed
system, differentiates the per-frame observable value with respect to the
parameter vector when needed, and then combines those per-frame values with the
reweighting correction that accounts for the change in state weights. The user
therefore only provides the observable itself; replay, Enzyme differentiation,
state reweighting, ESS tracking, normalization by tolerance, and line-search
acceptance are all handled by the generic optimizer.

\paragraph{Structured observables that need custom replay reduction.}

Some observables are not naturally represented as ``one scalar per frame and
then take a weighted mean.'' The dielectric constant is the current example:
its prediction depends on the weighted combination of
$\langle \mu^2 \rangle$, $\langle \mu \rangle$, and $\langle V \rangle$ rather
than on a single scalar frame observable. To keep the optimizer general, the
code exposes observable-specific dispatch hooks instead of hard-coding such
logic into the pipeline control flow.

Users who need a more structured observable can extend the following methods
for their own observable type:

\begin{itemize}
    \item \texttt{validate\_state\_observable\_target(observable, target\_name, leg\_cfg)} for target-level validation,
    \item \texttt{build\_state\_observable\_cache(observable, leg, params, param\_names, chain\_rule\_multiplier, state\_idx; compute\_gradients=true)} to define what per-frame quantities should be replayed and cached,
    \item \texttt{compute\_state\_observable\_estimate\_from\_cache(observable, leg, \\trainable\_param\_names, energy\_gradients\_phi, cache, energies\_current, log\_mixture\_denom, state\_idx, beta\_val, volumes; \\compute\_gradients=true, frame\_indices=nothing)} to reduce that cache into a prediction, gradient, and ESS,
    \item \texttt{state\_observable\_cache\_frame\_count(cache)} so that split-half confidence analysis can partition the cached observable data without knowing its internal representation.
\end{itemize}

In other words, a user who needs only a simple replayed scalar should supply a
callable and stop there, while a user who needs a composite ensemble observable
should extend these observable hooks for a custom observable type. The
surrounding machinery---target resolution, frozen-bias production, target loss
assembly, split-half confidence scaling, line search, and trust-region
moderation---remains unchanged.

\paragraph{Mixed-target optimization.}

Free energies and observables can be combined freely in one run. The optimizer
places them on a common footing by converting every target residual into
``tolerance units'' before applying the weighted Huber loss. Consequently, the
free-energy target may be in \(\mathrm{kcal/mol}\), a density target in
\(\mathrm{g/mL}\), a dielectric target in dimensionless relative permittivity,
and any user-defined observable in arbitrary units, provided each target is
given a meaningful tolerance. This normalization is what makes the pipeline a
general multi-target replay optimizer rather than a hydration-free-energy
special case.

\subsection{Offline Optimization and Step Acceptance}

The optimizer evaluates the current candidate parameter vector on the frozen production artifacts and assembles three objects:

\begin{enumerate}
    \item the full set of replayed target predictions,
    \item the summed tolerance-normalized objective and its latent-space gradient,
    \item the empirical Fisher matrix in latent space.
\end{enumerate}

The per-frame force-field gradients $\nabla_\theta U_\ell(\mathbf{x}_n, i; \theta)$ required for both the cycle gradient and the empirical Fisher estimate are obtained by differentiating the Molly potential-energy function with respect to the parameter vector via reverse-mode automatic differentiation, using Enzyme \citep{moses2020enzyme}. The chain-rule factor $d\theta_p/d\phi_p$ from the family-specific latent-to-physical transform is then applied to convert these $\theta$-space gradients into latent-space gradients before the Fisher and cycle-gradient assembly.

The base step is formed from the replay gradient and empirical Fisher matrix derived in the theory section, then stabilized through diagonal preconditioning, eigentruncation, KL-target scaling, and a per-pool infinity-norm cap on the latent update. The per-pool cap bounds the maximum absolute latent-coordinate change within each parameter pool independently, using a configurable \texttt{max\_phi\_step} per pool; this allows, for example, a tighter cap on Lennard-Jones $\sigma$ parameters than on partial-charge parameters. These are practical numerical design choices rather than additional formal identities. Before any update is attempted, the optimizer checks whether every target is already within its configured tolerance band; if so, that macro epoch exits immediately without taking another inner-loop step. The line search then proposes

\begin{equation}
    \phi_{\mathrm{prop}}
    =
    \phi_{\mathrm{current}}
    -
    \alpha \, \Delta \phi,
\end{equation}
%
starting from $\alpha=1$ and halving up to a fixed maximum of seven iterations until three conditions hold simultaneously:

\begin{itemize}
    \item the effective sample size, evaluated at the \emph{proposed} parameters $\phi_{\mathrm{prop}}$, remains above the per-leg threshold for every leg,
    \item the per-pool physical parameter drift relative to the reference model remains within configured bounds for every pool,
    \item the normalized objective satisfies the acceptance test described below.
\end{itemize}

The ESS is recomputed at $\phi_{\mathrm{prop}}$ on each halving step, not at the current parameters; it therefore checks whether the frozen reference ensemble still provides adequate importance-sampling support for the proposed candidate ensemble. The drift check is applied to the mapped physical parameters $\theta_{\mathrm{prop}} = \theta(\phi_{\mathrm{prop}})$ and gates per-pool families independently.

The objective test is not a simple monotonic decrease. Let $\tau$ be a configurable relative noise tolerance and let $r_{\mathrm{add}} \geq 0$ be the additional improvement requirement from split-half confidence (defined below). Denote by $\mathcal{J}$ the summed weighted Huber objective in normalized-residual space. A proposal is accepted when

\begin{equation}
    \mathcal{J}_{\mathrm{prop}}
    \leq
    \max\!\left(\varepsilon_{\mathrm{floor}},\;
        \mathcal{J}
        +
        \tau \, \mathcal{J}
        -
        r_{\mathrm{add}}
    \right),
    \label{eq:lsaccept}
\end{equation}
%
where $\varepsilon_{\mathrm{floor}}$ is a small positive floor derived from the noise tolerance. The window therefore widens by a fraction $\tau$ of the current objective magnitude, while $r_{\mathrm{add}}$ tightens it when production-data confidence is poor; the floor prevents the acceptance threshold from collapsing to zero. Using $\mathcal{J}$ instead of a raw-unit residual is what keeps mixed free-energy and observable targets on a common footing. If all seven halvings are exhausted without satisfying all three conditions the line search is declared failed, which triggers an early exit from the inner optimization loop.

\paragraph{Split-half confidence scaling.}
The production artifacts are also split deterministically into first and second temporal halves. The code recomputes target predictions, normalized target losses, and gradients on each half and forms three disagreement measures:

\begin{itemize}
    \item the maximum per-target prediction disagreement expressed in tolerance units,
    \item the disagreement in the full normalized objective,
    \item the disagreement in the objective gradient vector.
\end{itemize}

This is an implementation heuristic: split-half disagreement is used as a noise proxy for insufficiently reliable production data, not as a proof of correctness. It is conservative because it can catch obvious instability, but it can still miss structured long-timescale correlation. More formal block/bootstrap or replica-aware uncertainty estimators are natural future replacements rather than current behavior. These disagreement signals are compressed into a confidence factor

\begin{equation}
    s_{\mathrm{conf}}
    =
    \max
    \left(
        s_{\min},
        \frac{1}{1 + \kappa \, \delta_{\max}}
    \right),
\end{equation}
%
where $\delta_{\max}$ is the largest disagreement signal. The same split-half analysis also adds an extra objective-improvement requirement proportional to the normalized-objective disagreement before a line-search step is accepted.

\paragraph{Inner-loop exits.}

If the accepted line-search step becomes too small, if the line search repeatedly hits a tiny $\alpha$, or if the configured inner-epoch cap is reached, the code restores the best state seen inside that macro epoch and returns to fresh resimulation. A fresh macro-epoch Stage B failure does not trigger a hard rejection of the previously accepted parameter set; it simply prevents additional optimization until new readiness is established.


\subsection{Charge-Equilibration Parameterization}

When charge training is enabled, the code uses a constrained charge-equilibration model as a model-specific parameterization layer on top of the general replay-and-optimization machinery. Rather than optimizing partial charges $q_i$ directly, the trainable physical parameters are the electronegativity $\chi_i$ and hardness $\eta_i$ associated with each atom type or class.

For a molecule with target net charge $Q_{\mathrm{mol}}$, the partial charge on trainable atom $i$ is written as
\begin{equation}
    q_i
    =
    -\frac{\chi_i + \lambda_{\mathrm{mol}}}{\eta_i},
    \label{eq:qi_charge}
\end{equation}
%
where $\lambda_{\mathrm{mol}}$ is the molecule-specific Lagrange multiplier that enforces the net-charge constraint. Requiring
\begin{equation}
    \sum_{i \in \mathrm{trainable}} q_i
    =
    Q_{\mathrm{mol}} - \sum_{j \in \mathrm{fixed}} q_j,
\end{equation}
%
yields the analytic solution
\begin{equation}
    \lambda_{\mathrm{mol}}
    =
    -
    \frac{
        Q_{\mathrm{mol}} - \sum_{j \in \mathrm{fixed}} q_j
        +
        \sum_{i \in \mathrm{trainable}} \frac{\chi_i}{\eta_i}
    }{
        \sum_{i \in \mathrm{trainable}} \frac{1}{\eta_i}
    }.
    \label{eq:lagrange_mol}
\end{equation}

This construction enforces the molecular net charge exactly by construction while still allowing the optimizer to work in the bounded latent coordinates described earlier. In practice the hardness variables are kept positive through the chosen latent-to-physical transform so that the charge solve remains well conditioned, whereas electronegativity coordinates may be lightly bounded or left effectively unbounded depending on the training setup.

\subsection{Current Scope and Default Example}

The repository now supports a general thermodynamic cycle with arbitrary legs, coefficients, per-leg lambda schedules, and per-leg ensembles. The default example remains ethanol hydration. In that default setup:

\begin{itemize}
    \item the solvent leg uses an NPT expanded ensemble with a staged, charge-first and Lennard-Jones-second lambda schedule concentrated in the low-Lennard-Jones tail,
    \item the vacuum leg uses a simpler non-periodic leg with no pressure-volume term,
    \item the optimization can jointly train Lennard-Jones parameters and partial charges through the constrained charge-equilibration parameterization described above,
    \item the default target-mixing rule uses a $0.1\%$ tolerance band for both cycle free energies and replayed observables unless a target-specific tolerance is supplied explicitly,
    \item the default thermodynamic state is $20\,^\circ\mathrm{C}$ ($293.15\,\mathrm{K}$) so that the bundled solvent-density target is evaluated at the matching temperature,
    \item optimization is solute-only unless solvent parameters are explicitly enabled.
\end{itemize}

This is therefore no longer a speculative milestone roadmap. The implemented method is a readiness-gated, frozen-bias replay pipeline with support-aware Stage B validation, longer and thinner frozen production for noisy observables, tolerance-normalized multi-target optimization, bounded latent-coordinate updates, KL-limited natural-gradient steps, and deterministic confidence moderation from split-half production artifacts.
